% =====================================================================
\section{Introduction}
\label{sec:intro}
% =====================================================================

Software engineers are familiar with content-addressing.  Git identifies
every commit, tree, and blob by the SHA-1 hash of its
content~\cite{git}.  The Nix package manager derives store paths from
the hash of all build inputs~\cite{dolstra2004nix}.  IPFS addresses
files by their cryptographic hash, making retrieval
location-independent~\cite{benet2014ipfs}.  The Unison programming
language names functions by the hash of their syntax
tree~\cite{unison}.

In each case, the core idea is the same: derive the identifier from the
\emph{content itself}, not from a human-chosen name or a filesystem
location.  The result is an identifier that is \emph{stable} (it does
not change when the content moves), \emph{universal} (it requires no
central authority), and \emph{self-authenticating} (anyone can verify it
by rehashing).

\simone{Reduce Git analogy repetition---appears four times across Sections 1, 3, 4, 6. Once is enough.}

Formal mathematics has not yet adopted this pattern.  Theorems in proof
assistants---Lean~\cite{lean4}, Coq~\cite{coq}, Isabelle~\cite{isabelle},
Agda~\cite{agda}---are identified by fully qualified names tied to
project structure, file paths, and namespace conventions.  Renaming a
file, moving a declaration to a different module, or reorganizing a
namespace hierarchy breaks every downstream reference, even though the
mathematical content is unchanged.  This problem is not hypothetical: it
occurs routinely in Mathlib, the largest Lean~4 library, which undergoes
continuous refactoring across its ${\sim}200{,}000$
declarations~\cite{mathlib4}.  A complementary line of work addresses
the downstream consequence---automated proof
repair~\cite{ringer2021repair, gopinathan2023repair}---but does not
eliminate the root cause: names are not intrinsic to the mathematics.

This paper asks: \emph{can we apply content-addressing to formal
mathematics?}  We show that the answer is yes, with important caveats
about the granularity of equivalence (Section~\ref{sec:equiv}).  We
describe a working implementation for Lean~4 and a decentralized
registry architecture that requires no infrastructure beyond Git.

\paragraph{Contributions.} A content-addressing pipeline for Lean~4 declarations:
canonicalization, serialization, and SHA-256 hashing (Section~\ref{sec:pipeline}).
Two canonicalization levels with different stability/equivalence trade-offs (Section~\ref{sec:canon}).
A Merkle DAG hashing mode that gives declarations truly content-determined identity across project boundaries (Section~\ref{sec:modes}).
A decentralized registry architecture exploiting the monotonicity of mathematical proof (Section~\ref{sec:registry}). Empirical evaluation on 57{,}613 theorems from a Mathlib
module environment (Section~\ref{sec:results}). An open-source implementation as a Lean~4 library and CLI tool (Section~\ref{sec:impl}). A working implementation
is available at \href{https://github.com/marcellop71/CA}{https://github.com/marcellop71/CA}

\moqian{Prose needs significant rewriting. Currently reads as AI-generated. Rewrite for clarity, academic voice, narrative coherence.}

% --------------------------------------------------------------------
